\documentclass[11pt]{article}
\usepackage[a4paper,margin=2.5cm]{geometry}
\usepackage[T1]{fontenc}
\usepackage{lmodern}
\usepackage{amsmath,amssymb,amsthm}
\usepackage{booktabs,longtable,array}
\usepackage{enumitem}
\usepackage{listings}
\usepackage[hidelinks]{hyperref}
\newtheorem{theorem}{Theorem}[section]
\newtheorem{proposition}[theorem]{Proposition}
\theoremstyle{definition}
\newtheorem{definition}[theorem]{Definition}
\theoremstyle{remark}
\newtheorem{remark}[theorem]{Remark}
\lstset{basicstyle=\ttfamily\small,columns=fullflexible,keepspaces=true,breaklines=true,
  literate={₁}{{$_1$}}1 {₂}{{$_2$}}1 {ℓ}{{$\ell$}}1 {→}{{$\to$}}1 {∀}{{$\forall$}}1 {∃}{{$\exists$}}1 {≠}{{$\neq$}}1 {¬}{{$\neg$}}1 {↔}{{$\leftrightarrow$}}1 {∧}{{$\wedge$}}1}

\title{A typed bookkeeping language for verifiable research with language models:\\
\large design, Lean verification, and a case study}
\author{Marc Brendecke\\ \small Independent researcher, Bochum, Germany}
\date{Draft, \today}

\begin{document}
\maketitle

\begin{abstract}
Language models can carry out long research projects --- computing, formalizing, searching the literature and writing --- but over
hundreds of pages they produce characteristic errors that are hard to see in prose: quantities from different reference frames are
compared as if they were the same, values that were put in are later read as derived, and statements are promoted beyond what was
shown.  We describe QAM, a small typed language on top of Zermelo--Fraenkel set theory that makes these steps ill-formed, together with
an instruction layer (MML) and a store machine (MMA) whose ledger records every step.  The core is formalized in Lean~4 without
\texttt{sorry}: the eight object kinds are disjoint, a reindexing of the tag namespace is exactly invertible, admissible transport chains
preserve closure classes and compatible readouts, and the store machine satisfies preservation, progress for its non-modifying
instructions, atomicity, ledger completeness and the impossibility of a certificate for a false statement.  We report a case study: a
months-long project in mathematical physics, carried out by one author with language models from the start, with about a hundred
append-only audit rounds and more than sixty pre-registered tests.  We describe which classes of error the discipline exposed, which it
did not, and what the formal guarantees do and do not mean: the language secures the form of a record, not the truth of its content.
\end{abstract}

\section{Introduction}\label{sec:intro}

Language models can now take part in research well beyond answering questions.  They write and run numerical code, produce formal
proofs, search the literature and draft text, over sessions that add up to months.  For a single researcher without institutional
support this changes what is feasible.  It also changes what can go wrong.  A long record written in dialogue with a model accumulates
errors of a particular kind, and they are hard to see because each step reads plausibly in prose.  In the project described in
Section~\ref{sec:case} three such classes dominated:
\begin{enumerate}
\item \emph{Frame mixing.}  A quantity defined in one reference frame, or of one kind (a scale, a readout), is compared with a quantity
from another as if they were the same.
\item \emph{Input read as output.}  A value that was put in, or chosen to fit, is later cited as if it had been derived.
\item \emph{Promotion.}  A local, conditional or numerical statement is restated as a theorem, and later steps build on the stronger form.
\end{enumerate}
Peer review would catch some of these, but a long record produced at the speed of a model is not reviewed step by step, and the model
itself is not a reliable reviewer of its own earlier output.

This paper describes a pragmatic answer: a small typed language, QAM, in which the record is kept, and a store machine, MMA, which
executes instructions on the record and logs every step.  The language does not try to formalize the mathematics or the physics.  It
types the \emph{objects of the record} --- states, layers, operators, readouts, branches, shells, witnesses and tensor composites ---
and makes the three error classes ill-formed: a comparison across frames needs a declared comparison object, a value that rests on an
input carries a witness, and a certificate for a statement can only be formed from a proof.  The core of the language and of the machine
is formalized in Lean~4~\cite{lean4}.

\paragraph{Contributions.}
\begin{enumerate}
\item The language QAM over Zermelo--Fraenkel set theory, with its two-sided display for frame-typed quantities, and the instruction
layer MML with its transition semantics MMA (Sections~\ref{sec:language} and~\ref{sec:machine}).
\item Lean proofs of the properties that make the language usable as a record: disjointness of object kinds, an exactly invertible
reindexing, preservation of closure classes along admissible transport, and for the machine preservation, progress for the
non-modifying instructions, atomicity, ledger completeness and the absence of false certificates; and a project-independent module that
states the three error classes as rules on derivation trees and proves that well-formed trees never promote, never lose or invent an
input, and never change frame except through declared maps (Section~\ref{sec:lean}).
\item A \LaTeX{} package, \texttt{qam.sty}, that applies the same rules while a document is compiled and stops the compiler at the first
violation, that generates the QAM notation of each claim from its declaration, that links every claim to its sources by a chain of hashes starting at the thesis, and that exports
the record to a file which a verified Lean checker reads (Section~\ref{sec:tex}).
\item A case study of a research project in mathematical physics, carried out by one author with language models from the beginning,
with about a hundred append-only audit rounds and more than sixty pre-registered tests (Section~\ref{sec:case}).
\item An account of what the discipline caught and what it did not, and of what the formal guarantees mean (Section~\ref{sec:limits}).
\end{enumerate}

The guarantees are about form.  QAM can ensure that a record never compares across frames without saying so, never cites an input as a
derivation, and never certifies what was not proved.  It cannot ensure that a premise is true, that a literature value was read
correctly, or that a physical reading makes sense.  Section~\ref{sec:case} contains an example of each.

\section{The language}\label{sec:language}

\subsection{Object kinds}

QAM adds no axioms to Zermelo--Fraenkel set theory.  It fixes eight kinds of record objects as tagged tuples, the first entry being a
type tag:
\[
\begin{array}{llll}
\text{state} & \langle 0,a\rangle, &
\text{layer} & \langle 1,\ell,D\rangle,\\
\text{operator} & \langle 2,\ell_1,\ell_2,G\rangle, &
\text{readout} & \langle 3,\ell,W,R\rangle,\\
\text{branch} & \langle 4,b,B\rangle, &
\text{shell} & \langle 5,r,s,T\rangle,\\
\text{witness} & \langle 6,u\rangle, &
\text{tensor} & \langle 7,p,q\rangle .
\end{array}
\]
Here \(D\) is the carrier of a layer, \(G\) the graph of an operator from layer \(\ell_1\) to layer \(\ell_2\), \(W\) the set of readout values and \(R\)
the readout rule; a branch records a choice, a shell a refinement step, a witness an external input, and a tensor an admitted composite.
Recognition predicates such as \(\mathrm{QState}(x):\Leftrightarrow\exists a\,(x=\langle0,a\rangle)\) identify the kinds.  They are recognition
clauses, not an independent axiom system: anything that can be said in QAM can be said in ZF.

\subsection{Relations}

Three relations carry the content of a record: \(\mathrm{QAdmTransport}(o,s,t)\), the operator \(o\) carries the state \(s\) admissibly to
\(t\); \(\mathrm{QClEq}(s,t)\), \(s\) and \(t\) lie in the same closure class; and \(\mathrm{QRead}(r,s,w)\), the readout \(r\) reads \(s\) as the value \(w\).
Two schemata relate them: closure equivalence is an equivalence relation, and admissible transport preserves the closure class.  A readout
\(r\) is \emph{closure-compatible} for a map \(J\) on values if states in one closure class are read to values with the same image under \(J\).
What ``admissible'' and ``closure class'' mean is left to the application.  The schemata constrain them but do not fix them; both extremes ---
every transport admissible, or none --- satisfy them (Section~\ref{sec:lean}).

\subsection{Frames and the two-sided display}

Quantities that live in a frame carry it as an index.  In the two-sided display \({}^{A}{}_{B}\mathcal T^{i}{}_{j}\) the front indices record the
frame or readout side and the back indices the structure side.  Only an upper index may be contracted with a lower one, which is the ordinary
action or readout; two upper or two lower indices may be compared only through an explicitly declared comparison tensor, for example
\(G_{MT}\,{}^{M}X\,{}^{T}Y\) for a quantity \(X\) of frame \(M\) and \(Y\) of frame \(T\).  A comparison across frames without such an object is not
a well-formed expression.  This is the rule that addresses frame mixing.

\subsection{Witnesses, branches and status}

A value that rests on an input is written as a readout of a witness, \(w=R(u)\) with \(u=\langle6,\dots\rangle\), and a choice made during a
construction is a branch.  A derivation is therefore visibly different from a fit: the former has no witness in its premises.  This is the
rule that addresses input read as output.  Every statement of a record carries one of the status tags \emph{proved}, \emph{computed},
\emph{conditional}, \emph{reading}, \emph{open} or \emph{withdrawn}, and a statement can only be cited with its tag.

\subsection{Reindexing}

When the tag namespace is extended, existing tags are shifted by one and the tag \(0\) is reserved for ``no active kind''.  The shift and its
back-translation must be exact on inhabited tags, so that old records remain readable.

\section{The store machine}\label{sec:machine}

A record is kept in a \emph{store}: a content (objects at addresses and relations between addresses), a guard register (predicates the
content must satisfy), a list of snapshots, and a ledger.  A store is \emph{admissible} if every relation uses stored addresses and every
guard holds.

\paragraph{Instructions (MML).}  \textsc{store}\((a,x)\) puts an object at an address, \textsc{link}\((a,b)\) adds a relation,
\textsc{snapshot} saves the content, \textsc{rollback}\((k)\) restores the \(k\)-th snapshot, \textsc{scrub} is a no-op used for cleaning
passes, and \textsc{export}\((a)\) releases a stored object.

\paragraph{Transition semantics (MMA).}  Each instruction has a raw effect, which may be undefined (a rollback to a missing snapshot, an
export of a missing address).  A step executes the raw effect and checks the guards on the result:
\[
\mathrm{exec}(s,i)=\begin{cases}
(s\ \text{with entry } (i,\textsc{reject}),\ \textsc{reject}) & \text{if the raw effect is undefined},\\
(s'\ \text{with entry } (i,\textsc{cont}),\ \textsc{cont}) & \text{if it gives } s' \text{ and } s' \text{ is admissible},\\
(s\ \text{with entry } (i,\textsc{guard}),\ \textsc{guard}) & \text{otherwise.}
\end{cases}
\]
A result that violates a guard is never committed; the old content stays and the ledger records the defect.  A program is a list of
instructions executed in order.

\paragraph{Certificates.}  The record distinguishes a \emph{candidate}, which the machine can produce for any statement, from a
\emph{certificate}, which contains a proof.  A candidate becomes a certificate only through a proof supplied from outside the machine.  No
instruction produces a certificate.

\section{What is proved in Lean}\label{sec:lean}

The language and the machine are formalized in Lean~4 (version 4.35.0) in four modules of the library \texttt{QSFormal}~\cite{repo}.  They
use the Lean core only; other modules of the same library, on the physics, use Mathlib~\cite{mathlib}.  Continuous integration builds the
library and rejects \texttt{sorry}, \texttt{admit} and \texttt{native\_decide}.  Theorem names below are the Lean names.

\subsection{Object kinds and reindexing}

The eight kinds are the constructors of an inductive type \texttt{Obj}; the ZF tag is a function \texttt{Obj.tag}.

\begin{theorem}[\texttt{qState\_iff\_tag}, \texttt{not\_state\_and\_op}, \texttt{tag\_lt\_eight}]
An object is a state exactly when its tag is \(0\); no object is both a state and an operator; every tag is less than eight.
\end{theorem}

\begin{theorem}[\texttt{back\_shift}, \texttt{shift\_back}, \texttt{back\_toShifted}, \texttt{toShifted\_back}]
The shift \(k\mapsto k+1\) and the back-translation \(a\mapsto a-1\) on inhabited tags \(a\ge1\) are mutually inverse, and the shifted namespace
with reserved \(0\) is in bijection with ``no kind or a kind''.
\end{theorem}

\subsection{Transport and readout}

A \emph{model} of the transport layer is a structure with relations \texttt{admTransport}, \texttt{clEq}, \texttt{read} and the two schemata as
fields.  Every theorem holds for every model.

\begin{theorem}[\texttt{first\_transport\_principle}, \texttt{chain\_clEq}, \texttt{chain\_readout}]
An admissible transport keeps a state in its closure class, and a closure-compatible readout gives the same value, up to \(J\), before and
after.  The same holds for every finite chain of admissible transports.
\end{theorem}

\begin{theorem}[\texttt{totalModel}, \texttt{emptyModel}, \texttt{empty\_chain}]
The schemata have two trivial models: in one every transport is admissible and there is one closure class, in the other no transport is
admissible and the classes are singletons.  In the second, a chain of admissible transports is trivial.
\end{theorem}

The last theorem is a warning as much as a result.  The schemata are consistent, but they do not determine what admissible means; all
content enters through the choice of the model.

\paragraph{Routes.}  For frame-typed records there is one model per route \(a\to b\), with a composition law as an explicit assumption: an
admissible transport from \(a\) to \(b\) followed by one from \(b\) to \(c\) is admissible from \(a\) to \(c\).

\begin{theorem}[\texttt{path\_admissible}, \texttt{path\_clEq}, \texttt{loop\_clEq}]
A path through any number of frames is admissible on the direct route and stays in a closure class of it; in particular a loop
\(a\to\dots\to a\) never changes the class of the route \(a\to a\).
\end{theorem}

\subsection{The store machine}

\begin{theorem}[Preservation: \texttt{exec\_preserves}, \texttt{run\_preserves}, \texttt{exec\_cont\_admissible}]
From an admissible store no instruction, and no program, leads to an inadmissible one; a step with status \textsc{cont} ends in an admissible
store.
\end{theorem}

\begin{theorem}[No stuck states and exact status: \texttt{exec\_ne\_halt}, \texttt{exec\_cont\_iff}, \texttt{exec\_reject\_iff}]
No step ends with status \textsc{halt}.  A step has status \textsc{cont} exactly when its raw effect is defined and admissible, and status
\textsc{reject} exactly when its raw effect is undefined.
\end{theorem}

\begin{theorem}[Atomicity: \texttt{exec\_atomic}]
A step that does not end with \textsc{cont} leaves content, snapshots and guards unchanged.
\end{theorem}

\begin{theorem}[Progress: \texttt{snapshot\_progress}, \texttt{scrub\_progress}, \texttt{export\_progress}]
In an admissible store, \textsc{snapshot} and \textsc{scrub} always succeed, and \textsc{export}\((a)\) succeeds for every stored address \(a\).
\end{theorem}

\begin{theorem}[Lineage: \texttt{exec\_logs}, \texttt{exec\_ledger\_length}, \texttt{run\_ledger\_length}]
Every step puts an entry with its instruction and status at the head of the ledger, and a program adds exactly as many entries as it has
instructions.
\end{theorem}

\begin{theorem}[Rollback and certificates: \texttt{rollback\_restores}, \texttt{no\_false\_certificate}]
A rollback to an existing snapshot whose content is admissible restores it exactly.  There is no certificate for \texttt{False}, whatever
instructions were run.
\end{theorem}

The statements follow the standard shape of type soundness~\cite{wright,pierce} --- preservation and progress --- adapted to a machine whose
steps may be refused.  Progress is proved for the three instructions that do not change an admissible content; for \textsc{store},
\textsc{link} and \textsc{rollback} it depends on the guards by design, and the exact status characterization says what happens instead.
The certificate rule is a small instance of the idea behind proof-carrying code~\cite{necula}: what is certified carries its proof.

\subsection{Status, inputs and frames}

The three error classes of Section~\ref{sec:intro} are stated directly in a module \texttt{QAM.Provenance} of a separate Lean package
\texttt{qam/lean}, which uses the Lean core only (no Mathlib) and does not depend on the rest of the library, or on any physics.  A \emph{claim} has a status, a frame and a list of inputs, the identifiers of the values it rests
on.  Statuses are ordered, \(\textsf{open}<\textsf{reading}<\textsf{conditional}<\textsf{computed}<\textsf{proved}\).  Claims are combined by
derivation trees with three kinds of node: a \emph{step} derives a statement from one premise, a \emph{join} from two premises, and a
\emph{conversion} moves a statement into another frame through a map \(g\).  A tree is well formed relative to a set \(D\) of declared maps
when no node has a higher status than its premises, a join combines premises of one frame only, and every conversion uses a map in \(D\)
with the right source frame.

\begin{theorem}[No promotion: \texttt{no\_promotion}, \texttt{proved\_only\_from\_proved}]
In a well-formed tree the conclusion has at most the status of every premise.  In particular a conclusion with status \textsf{proved}
rests only on proved claims.
\end{theorem}

\begin{theorem}[Inputs: \texttt{inputs\_tracked}, \texttt{inputs\_sourced}, \texttt{independent\_iff}]
The inputs of a conclusion are exactly the inputs of its premises: none is lost and none is invented.  A conclusion is independent of an
input exactly when no premise uses it.
\end{theorem}

\begin{theorem}[Frames: \texttt{frames\_declared}, \texttt{no\_maps\_same\_frame}]
In a well-formed tree every premise reaches the frame of the conclusion through a chain of declared maps.  Without declared maps, all
premises live in the frame of the conclusion.
\end{theorem}

These theorems are elementary, and they are meant to be: the point is that the rules are precise enough to be checked, by a proof
assistant and, as the next section shows, by the \TeX{} compiler.

\paragraph{Simplifications.}  Addresses are natural numbers, relations carry no pointer conditions, and the instructions that clone or merge
stores are not formalized.

\section{Case study}\label{sec:case}

\subsection{Setting}

The project is a programme in mathematical physics that starts from one conjecture, by the author: that matter, light and a tachyon
sector are one object read out in three frames~\cite{companion,physics}.  Everything else grew out of it in a dialogue with language models
(Claude; ChatGPT and Codex) used from the beginning.  By November 2026 the research text had grown to about \(1300\) pages, of which roughly
\(85\%\) are protocols, registers and audit rounds.  The author is not a physicist; the models wrote most code, the Lean formalization and much
of the text.

\subsection{Process}

\begin{itemize}
\item \emph{Append-only rounds.}  Results are added as numbered rounds; earlier rounds are never rewritten, and a correction is a new round
that names what it corrects.  Since the first audited release there have been about a hundred such rounds.
\item \emph{Pre-registration.}  Every test with a decision criterion is written down before the computation --- hypothesis, data, criterion ---
and committed; the commit time stamps it~\cite{nosek}.  More than sixty tests have been registered.  Notes that were not blind say so.
\item \emph{Scope list and linter.}  A fixed list of statements that are not claimed is checked against every round, also mechanically.
\item \emph{Recomputation.}  Results that carry weight are recomputed with independently written code, at higher precision where needed,
and in one series cross-checked between different models.  Large rank computations run in parallel on continuous-integration runners.
\item \emph{Review.}  Before a paper is released, a separate model instance without access to the record reads only the manuscript and writes
a referee report.
\end{itemize}

\subsection{What the register shows}

Of the registered tests, the register lists more than a dozen as refuted, six as confirmed, and six as invalid by their own criterion ---
for instance because a required numerical gap was not reached, even though the numbers came out as predicted.  Invalid tests are recorded as
invalid, not as passed.  Four predictions wait for new data.  Among the statements of the research text that were withdrawn or corrected are
all steps that relied on an external claimed proof of the Riemann hypothesis, an argument for the cosmological constant that is off by a factor
of about \(10^{121}\), and six early predictions, none of which held.

\subsection{Error classes and how they were caught}

\begin{center}
\small
\renewcommand{\arraystretch}{1.3}
\begin{longtable}{@{}>{\raggedright\arraybackslash}p{0.20\textwidth}>{\raggedright\arraybackslash}p{0.43\textwidth}>{\raggedright\arraybackslash}p{0.29\textwidth}@{}}
\toprule
class & example & caught by\\
\midrule
\endhead
frame mixing & A ``structural value'' of the central parameter equated a readout of one frame with an offset measured in another; a
second one labelled two different states with the same symbol & two-sided display: the comparison has no declared comparison tensor\\
input read as output & The central parameter, put in from two literature values, was cited as derived; a formula reproducing it to
\(0.018\%\) had an exponent chosen after the fact & witness typing: the value is a readout of a witness\\
promotion & Local or conditional statements stated as theorems and built upon, e.g.\ a cosmological-constant chain & status tags and the scope
list; audit rounds\\
external dependence & Proofs that cited a separate paper claiming the Riemann hypothesis & \emph{not} by the language: an audit round traced
the citations\\
numerical artefact & A high-precision result that came from passing 100-digit numbers through a 15-digit text conversion & recomputation at a
generic point; withdrawn before it entered a round\\
literature reading & A counting value read as ``corrected'' that later literature considers less well founded & independent literature check\\
\bottomrule
\end{longtable}
\end{center}

\subsection{A translation table}

For eight key statements of the project the following table gives the standard notation and the QAM form.  The QAM form is not shorter;
it records what the standard form leaves to the surrounding text.

\begin{center}
\small
\renewcommand{\arraystretch}{1.3}
\begin{longtable}{@{}>{\raggedright\arraybackslash}p{0.27\textwidth}>{\raggedright\arraybackslash}p{0.37\textwidth}>{\raggedright\arraybackslash}p{0.28\textwidth}@{}}
\toprule
standard notation & QAM notation & what the QAM form adds\\
\midrule
\endhead
kernel \(q=it\,\mathbf 1+z_ke_k\in\mathbb C\otimes\mathbb H\) & \(\ell_n=\langle1,n,\mathbb H\otimes\mathbb C\rangle\), \(s_q=\langle0,q\rangle\), \(q\in D(\ell_n)\) & the layer on which the
state lives\\
compensator \(C=\tfrac12(1+J)\), \(C^2=C\) & \(o_C=\langle2,\ell,\ell,G_C\rangle\), \(G_C\circ G_C=G_C\) & an operator within one layer, not a readout\\
sector by the sign of \(\sigma=z\cdot z\) & \(r_\sigma=\langle3,\ell,\{+,0,-\},R_\sigma\rangle\) & the sector is a readout of the state\\
exchange of sectors \(z\mapsto i\bar z\) & \({}^{T}{}_{M}\mathcal J^{i}{}_{j}z^j=(i\bar z)^i\); comparison only via \(G_{MT}\) & source and target frame; raw
comparison ill-formed\\
compensator matrix, \([\mathcal M,J]=[\mathcal M,C]=0\) & \(G_{\mathcal M}\circ G_J=G_J\circ G_{\mathcal M}\), \(G_{\mathcal M}\circ G_C=G_C\circ G_{\mathcal M}\) & nothing beyond
the standard form\\
one generation from \(\mathbb C\otimes\mathbb H\otimes\mathbb O\) & \(\langle7,\mathbb C\otimes\mathbb H,\mathbb O\rangle\); the choice of \(J\) is a branch \(\langle4,b_J,B\rangle\) & the construction is
visibly conditional\\
frame offset \(\beta\approx0.549\) & \({}^{M}\beta=R_\beta(u)\), witness \(u=\langle6,(\gamma_1,\gamma_2)\rangle\) & the value rests on an input\\
matter fraction \(\Omega_M=\beta^2\) & \(\mathrm{QRead}(r_\Omega,s_{\rm now},\beta^2)\), layer \(\ell_{\rm now}\) & validity restricted to one epoch\\
\bottomrule
\end{longtable}
\end{center}

\section{Related work}\label{sec:related}

\paragraph{Proof assistants.}  Lean~\cite{lean4} and its library Mathlib~\cite{mathlib} make it possible to verify mathematics, and language
models are increasingly used to produce such proofs.  QAM is complementary: it does not verify the mathematics of a record, only its
bookkeeping, and only a small part of the project's mathematics is formalized.

\paragraph{Type soundness and proof-carrying code.}  The properties proved for the store machine are the familiar ones of syntactic type
soundness~\cite{wright,pierce}.  The rule that a certificate carries its proof echoes proof-carrying code~\cite{necula}.

\paragraph{Provenance.}  The ledger of the store machine is a provenance record in the sense of the W3C PROV model~\cite{prov}, reduced to
instructions and outcomes.  QAM adds typing of the recorded objects, which provenance models leave open.

\paragraph{Pre-registration.}  Writing down hypotheses, data and criteria before the analysis is established practice in parts of the
empirical sciences, as a response to problems of reproducibility~\cite{nosek,baker}.  The case study applies it to computational and
theoretical work, where a commit time stamp is enough to fix a test in advance.

\paragraph{Language models in research.}  Systems such as FunSearch~\cite{funsearch} search for mathematical constructions with language models
under automatic evaluation, and proposals for automated research pipelines exist~\cite{aiscientist}.  The present setting is different: a
single human directs a long, open-ended project, and the problem is not search but keeping an honest record over months.  We are not aware of a
documented case of comparable length in which the record itself is typed and partly verified; a systematic survey remains to be done.

\section{Limits and discussion}\label{sec:limits}

\paragraph{What the language does not prevent.}  QAM secures form, not truth.  It did not prevent the dependence on a false external proof; an
audit round found it.  It does not check whether a literature value is read correctly, whether a premise holds, or whether a physical reading is
sensible.  Its schemata admit trivial models, so a record can be well-typed and empty.  And a well-typed record can still be wrong in every
statement that is not about bookkeeping.

\paragraph{Cost.}  The discipline is expensive.  The research text is dominated by protocols, and a reader who wants the results needs a separate,
short account.  Status tags and witness typing slow down writing.  In the case study the cost was worth paying, because the alternative was a long
record that nobody, including the author, could trust; whether it is worth paying in other settings is open.

\paragraph{What helped most.}  In order of effect: pre-registration with time stamps, append-only rounds, recomputation with independent code,
the review of manuscripts by a model instance without access to the record, and the typing rules.  The typing rules were most useful for the class
of error they were designed for, frame mixing, and least useful for errors of premise.

\paragraph{Next steps.}  Progress for \textsc{store}, \textsc{link} and \textsc{rollback} under explicit guard conditions; the clone and merge
instructions; a converter that adds \texttt{qam.sty} declarations to an existing research text, so that the compiler checks it; and a comparison with electronic lab-notebook and
provenance systems.

\section{Conclusion}\label{sec:conclusion}

A small typed language over ordinary set theory, a store machine with a complete ledger, and a handful of Lean theorems about them are enough to
make three frequent errors of long, model-assisted research ill-formed, and to make every step of a record traceable.  They are not enough to
make the record true.  The case study suggests that the combination with pre-registration, append-only correction and independent review is what
makes such a record trustworthy, and that the typing earns its place in that combination as a guard against one specific, common and otherwise
hard to see kind of error.


\section*{Data availability}
The Lean sources, the research text, the audit rounds and the registration notes are in the public repository~\cite{repo}.

\appendix
\section{Lean modules and use of language models}\label{app:build}

\paragraph{Modules.}  \texttt{QSFormal/QAM.lean} (object kinds, reindexing, transport models, chains), \texttt{QSFormal/Models.lean} (the two
trivial models), \texttt{QSFormal/Route.lean} (routed models and paths), \texttt{QSFormal/Store.lean} (store, instructions, transition
semantics, preservation, rollback, certificates) and \texttt{QSFormal/StoreProgress.lean} (no stuck states, status characterization, atomicity,
progress, ledger completeness), built with Mathlib in the directory \texttt{formal/}.  The project-independent part is a separate
package in \texttt{qam/lean/} that uses the Lean core only: \texttt{QAM/Provenance.lean} (derivation trees: no promotion, inputs, frames)
and \texttt{QAM/QamFile.lean} (the checker for exported records, its line-by-line and record-level soundness, the hash links), with the
executable \texttt{qamcheck}.  Both are built with \texttt{lake build}; the toolchain is fixed in \texttt{lean-toolchain}.

\paragraph{Use of language models.}  The starting conjecture of the case study is the author's.  The language, the machine, the Lean
formalization and much of this text were developed in a dialogue with large language models (Claude, Anthropic; ChatGPT and Codex, OpenAI).
Some ideas came from the models, others from the author; the author chose the questions and decided what to keep.  Their output was accepted
only through the checks described in Section~\ref{sec:case}.  The author takes full responsibility for the content.


\begin{thebibliography}{99}
\bibitem{repo} M.~Brendecke, QuantumSphere-Unfold, \url{https://github.com/Moriarty1982/QuantumSphere-Unfold}.
\bibitem{companion} M.~Brendecke, \emph{Quantum Sphaera Companion: A Structure-First Mathematical Unfolding}, Zenodo, concept DOI
\href{https://doi.org/10.5281/zenodo.19210728}{10.5281/zenodo.19210728}.
\bibitem{lean4} L.~de Moura and S.~Ullrich, The Lean~4 theorem prover and programming language, in \emph{Automated Deduction --- CADE 28},
Lecture Notes in Computer Science \textbf{12699}, Springer 2021, 625--635.
\bibitem{mathlib} The mathlib Community, The Lean mathematical library, in \emph{Proceedings of the 9th ACM SIGPLAN International Conference
on Certified Programs and Proofs (CPP 2020)}, ACM 2020, 367--381.
\bibitem{wright} A.~K.~Wright and M.~Felleisen, A syntactic approach to type soundness, \emph{Information and Computation} \textbf{115}
(1994) 38--94.
\bibitem{pierce} B.~C.~Pierce, \emph{Types and Programming Languages}, MIT Press 2002.
\bibitem{necula} G.~C.~Necula, Proof-carrying code, in \emph{Proceedings of the 24th ACM Symposium on Principles of Programming
Languages (POPL 1997)}, ACM 1997, 106--119.
\bibitem{prov} W3C, PROV-DM: The PROV Data Model, W3C Recommendation, 30 April 2013.
\bibitem{nosek} B.~A.~Nosek, C.~R.~Ebersole, A.~C.~DeHaven and D.~T.~Mellor, The preregistration revolution, \emph{Proc.\ Natl.\ Acad.\ Sci.\ USA}
\textbf{115} (2018) 2600--2606.
\bibitem{baker} M.~Baker, 1,500 scientists lift the lid on reproducibility, \emph{Nature} \textbf{533} (2016) 452--454.
\bibitem{funsearch} B.~Romera-Paredes et al., Mathematical discoveries from program search with large language models, \emph{Nature}
\textbf{625} (2024) 468--475.
\bibitem{aiscientist} C.~Lu, C.~Lu, R.~T.~Lange, J.~Foerster, J.~Clune and D.~Ha, The AI Scientist: towards fully automated open-ended
scientific discovery, arXiv:2408.06292 (2024).
\bibitem{physics} M.~Brendecke, Three readouts of one kernel: duality sectors, time reversal and a compensator-form projector in
\(\mathbb C\otimes\mathbb H\otimes\mathbb O\), draft (2026), in~\cite{repo}.
\end{thebibliography}
\end{document}
